\documentclass{article}
\usepackage{amsmath}
\usepackage{amssymb}
\usepackage{amsfonts}
\usepackage{booktabs}
\usepackage{geometry}
\usepackage{tabularx}
\usepackage{array}
\geometry{margin=1in}

\title{STAT452: Variable Selection and Regularization --- Lesson 3 Notes}
\author{Nguyen Dinh Thien Loc}
\date{\today}

\begin{document}

\maketitle

\section{Motivation for Variable Selection}

\subsection{Why Perform Variable Selection?}
\begin{itemize}
    \item \textbf{Simplicity (Occam's Razor):} The simplest adequate model is best. Redundant predictors should be removed to maintain interpretability.
    \item \textbf{Noise Reduction:} Unnecessary predictors add noise and obscure true relationships, wasting degrees of freedom and reducing estimation precision.
    \item \textbf{Collinearity Mitigation:} Too many variables may attempt to explain the same effect, leading to unstable and inflated coefficient estimates.
    \item \textbf{Cost Minimization:} Fewer predictors mean lower data collection costs and faster models, which is especially important in predictive modeling.
\end{itemize}

\subsection{Preparation Before Selection}
\begin{itemize}
    \item Screen for \textbf{outliers and influential points}; consider removing or treating them temporarily.
    \item Apply \textbf{appropriate transformations} (normalization, log-transform, standardization) as needed.
\end{itemize}

\section{The Problem with $p$-Value-Based Selection}

\subsection{The Wald Test Statistic}
To test $H_0: \beta_i = 0$, we use the Wald test statistic:
$$t = \frac{\hat{\beta}_i}{\text{SE}(\hat{\beta}_i)}$$
Under $H_0$, this follows a $t_{n-(p+1)}$ distribution and tends to the standard normal as $n \to \infty$.

\subsection{Decomposition of the Test Statistic}
The test statistic can be decomposed as:
$$t = \frac{\hat{\beta}_i}{\displaystyle \frac{\hat{\sigma}}{\sqrt{n \cdot \widehat{\text{Var}}(X_i)}} \cdot \sqrt{\text{VIF}_i}}$$
The denominator (standard error) depends on:
\begin{enumerate}
    \item \textbf{Noise level ($\hat{\sigma}$):} More noise inflates the standard error $\Rightarrow$ lower significance.
    \item \textbf{Sample size ($n$):} Larger $n$ shrinks the standard error ($\sim 1/\sqrt{n}$) $\Rightarrow$ higher significance.
    \item \textbf{Predictor variance ($\widehat{\text{Var}}(X_i)$):} More spread in $X_i$ improves estimation $\Rightarrow$ higher significance.
    \item \textbf{Collinearity ($\text{VIF}_i$):} High multicollinearity inflates the standard error $\Rightarrow$ lower significance.
\end{enumerate}

\subsection{Why $p$-Values Are Misleading for Selection}
\begin{itemize}
    \item The test statistic reflects \textbf{data quality and predictor structure}, not just effect size.
    \item \textbf{Unimportant} variables with high variance or low multicollinearity may appear significant.
    \item \textbf{Important} variables may appear non-significant if measured with high uncertainty.
    \item As $n \to \infty$, even \textbf{tiny nonzero coefficients} will become statistically significant.
    \item $p$-value-based selection favors predictors with high variance, low correlation to others (low VIF), and large sample sizes --- regardless of practical relevance.
\end{itemize}

\noindent\textbf{Better alternatives:} Cross-validation (CV), Information criteria (AIC, BIC), Regularization methods (LASSO, Ridge).

\newpage
\section{Cross-Validation for Variable Selection}

Cross-validation estimates the model's ability to predict \textbf{unseen data}. While $C_p$, AIC, and BIC attempt this analytically, CV does it empirically.

\subsection{Leave-One-Out Cross-Validation (LOOCV)}
\begin{itemize}
    \item For each observation $i = 1, \dots, n$: fit the model on all data \textit{except} observation $i$, then predict $\hat{y}_{(i)}$.
    \item The LOOCV error is:
    $$\text{CV}_{(n)} = \frac{1}{n} \sum_{i=1}^{n} (y_i - \hat{y}_{(i)})^2$$
    \item \textbf{Properties:} Low bias (nearly the full dataset is used for training), but higher variance and computationally expensive ($n$ models to fit).
    \item Tends to \textbf{overfit} (includes unnecessary variables).
\end{itemize}

\subsection{$k$-Fold Cross-Validation}
\begin{itemize}
    \item Randomly partition the data into $k$ roughly equal-sized folds (typically $k = 5$ or $10$).
    \item For each fold $j$: train on the remaining $k-1$ folds, predict on fold $j$, and compute $\text{MSE}_j$.
    \item The $k$-fold CV error is:
    $$\text{CV}_{(k)} = \frac{1}{k} \sum_{j=1}^{k} \text{MSE}_j$$
    \item \textbf{Properties:} More computationally efficient (only $k$ models to train), less variance in error estimates, slightly less accurate if all models are misspecified.
\end{itemize}

\section{Information Criteria}

\subsection{Akaike Information Criterion (AIC)}
\begin{itemize}
    \item \textbf{Definition:}
    $$\text{AIC} = -2 \log \hat{L} + 2p$$
    where $\hat{L}$ is the maximum likelihood of the fitted model and $p$ is the number of parameters (including intercept).
    \item \textbf{Interpretation:} AIC estimates the relative out-of-sample prediction error. It penalizes model complexity to avoid overfitting. \textbf{Lower AIC values} indicate a better trade-off between fit and simplicity.
    \item \textbf{Use when:} Your goal is \textbf{prediction}, not necessarily finding the ``true'' model.
\end{itemize}

\subsection{Bayesian Information Criterion (BIC)}
\begin{itemize}
    \item \textbf{Definition:}
    $$\text{BIC} = -2 \log \hat{L} + p \log n$$
    where $n$ is the sample size.
    \item \textbf{Interpretation:} BIC penalizes complexity \textbf{more strongly} than AIC (because of the $\log n$ term). As $n \to \infty$, BIC tends to select the correct (true) model if it exists. \textbf{Lower BIC is better.}
    \item \textbf{Use when:} You prioritize \textbf{simpler models} and model identification over pure predictive performance.
\end{itemize}

\subsection{Mallows' $C_p$ Criterion}
\begin{itemize}
    \item \textbf{Definition:}
    $$C_p = \frac{\text{RSS}_p}{\hat{\sigma}^2} - n + 2p$$
    where $\text{RSS}_p$ is the Residual Sum of Squares from a model with $p$ predictors, $\hat{\sigma}^2$ is the error variance estimate from the \textbf{full model} (with all predictors), and $n$ is the number of observations.
    \item \textbf{Interpretation:} If the subset model is unbiased, then $C_p \approx p$. Models with $C_p \approx p$ and small overall value are preferred.
    \item \textbf{Use when:} You trust the full model as a benchmark to estimate $\hat{\sigma}^2$.
\end{itemize}

\subsection{Comparison: AIC vs. BIC vs. $C_p$}
\begin{table}[h]
\centering
\small
\begin{tabularx}{\textwidth}{>{\raggedright\arraybackslash}p{0.18\textwidth}
                            >{\raggedright\arraybackslash}p{0.30\textwidth}
                            >{\raggedright\arraybackslash}X}
\toprule
\textbf{Criterion} & \textbf{Penalty Strength} & \textbf{Best Use Case} \\ 
\midrule
AIC & Lenient ($2p$) & When \textbf{prediction} is the primary goal. \\
BIC & Stronger ($p \log n$) & When identifying the ``true'' model or a \textbf{simpler model} is preferred. \\
$C_p$ & Error-based & When the \textbf{full model} variance estimate $\hat{\sigma}^2$ is trustworthy. \\
\bottomrule
\end{tabularx}
\end{table}

\newpage
\section{Stepwise Selection Methods}

Stepwise methods add or remove one variable at a time based on a selection criterion (AIC, BIC, Adjusted $R^2$, Mallows' $C_p$, CV error, etc.).

\subsection{Forward Stepwise Selection}
\begin{enumerate}
    \item Start with \textbf{no predictors} (intercept-only model).
    \item At each step, add the variable that \textbf{most improves} the model (e.g., lowest AIC).
    \item Repeat until no further significant improvement is obtained.
\end{enumerate}

\subsection{Backward Stepwise Selection}
\begin{enumerate}
    \item Start with the \textbf{full model} (all predictors included).
    \item At each step, remove the variable whose exclusion \textbf{improves performance most}.
    \item Continue until further removal worsens the model.
\end{enumerate}

\subsection{Mixed (Forward-Backward) Stepwise Regression}
\begin{enumerate}
    \item Start with an initial model (e.g., intercept-only).
    \item Add the predictor that most improves the model (e.g., lowest AIC or highest $F$-statistic).
    \item After adding, \textbf{re-evaluate all included variables}.
    \item Remove any variable that no longer contributes significantly.
    \item Repeat steps 2--4 until no further change is possible.
\end{enumerate}
\noindent\textbf{Key point:} Variables already added may still be removed later --- this is the distinguishing feature from pure forward selection.

\subsection{Stepwise vs. All-Subsets Selection}
\begin{itemize}
    \item Stepwise methods only explore a \textbf{subset of possible models} and are computationally efficient but may miss the globally best model.
    \item All-subsets selection considers all $2^p$ models (feasible only for small $p$).
    \item Stepwise methods are a \textbf{compromise}: faster but potentially suboptimal.
\end{itemize}

\newpage
\section{Ridge Regression (L2 Regularization)}

\subsection{Limitations of OLS Motivating Regularization}
\begin{itemize}
    \item \textbf{Multicollinearity:} When predictors are highly correlated, $\mathbf{X}^T\mathbf{X}$ becomes nearly singular, producing unstable and highly sensitive coefficient estimates.
    \item \textbf{High Variance:} As $p$ approaches or exceeds $n$, OLS predictions become highly variable with poor generalization.
    \item \textbf{No Variable Selection:} OLS includes all predictors regardless of relevance, reducing interpretability.
\end{itemize}

\subsection{Ridge Regression Formulation}
Ridge regression solves the following penalized optimization problem:
$$\hat{\boldsymbol{\beta}}_{\text{ridge}} = \arg\min_{\boldsymbol{\beta}} \left\{ \sum_{i=1}^{n} \left( y_i - \beta_0 - \sum_{j=1}^{p} \beta_j x_{ij} \right)^2 + \lambda \sum_{j=1}^{p} \beta_j^2 \right\}$$

\noindent\textbf{Equivalent constrained form:}
$$\min_{\boldsymbol{\beta}} \sum_{i=1}^{n} \left( y_i - \beta_0 - \sum_{j=1}^{p} \beta_j x_{ij} \right)^2 \quad \text{subject to} \quad \sum_{j=1}^{p} \beta_j^2 \le t$$

\noindent\textbf{Matrix notation and closed-form solution:}
$$\hat{\boldsymbol{\beta}}_{\text{ridge}} = \arg\min_{\boldsymbol{\beta}} \left\{ \frac{1}{2} \|\mathbf{y} - \mathbf{X}\boldsymbol{\beta}\|^2 + \lambda \|\boldsymbol{\beta}\|^2 \right\} = (\mathbf{X}^T\mathbf{X} + \lambda \mathbf{I})^{-1} \mathbf{X}^T \mathbf{y}$$

\noindent\textbf{Key properties:}
\begin{itemize}
    \item Adding $\lambda \mathbf{I}$ to $\mathbf{X}^T\mathbf{X}$ ensures invertibility even under multicollinearity.
    \item The intercept $\beta_0$ is \textbf{not penalized}.
    \item Predictors should be \textbf{standardized} before fitting to ensure fair penalization across variables.
\end{itemize}

\subsection{Standardization and Intercept Treatment}
Before estimating Ridge coefficients:
\begin{itemize}
    \item \textbf{Intercept:} Set $\hat{\beta}_0 = \bar{y}$; the penalty applies only to slope coefficients $(\beta_1, \beta_2, \dots, \beta_p)$.
    \item \textbf{Standardization:} Transform predictors to ensure comparability:
    $$\tilde{x}_{ij} = \frac{x_{ij}}{\sqrt{\frac{1}{n} \sum_{k=1}^{n} (x_{kj} - \bar{x}_j)^2}}$$
    This prevents penalization from disproportionately affecting variables with larger scales.
\end{itemize}

\subsection{The Bias-Variance Trade-Off}
\begin{itemize}
    \item The Mean Squared Error decomposes as: $\text{MSE} = \text{Bias}^2 + \text{Variance}$.
    \item Increasing model complexity reduces bias but increases variance.
    \item Ridge slightly \textbf{increases bias} but \textbf{significantly reduces variance}, finding a favorable balance.
\end{itemize}

\subsection{Bias of Ridge Estimates}
\begin{itemize}
    \item \textbf{OLS estimate:} $\hat{\boldsymbol{\beta}}_{\text{OLS}} = (\mathbf{X}^T\mathbf{X})^{-1}\mathbf{X}^T\mathbf{Y}$
    \item \textbf{Ridge estimate:} $\hat{\boldsymbol{\beta}}_{\text{ridge}} = (\mathbf{X}^T\mathbf{X} + \lambda \mathbf{I}_p)^{-1}\mathbf{X}^T\mathbf{Y}$
    \item \textbf{Expected value:}
    $$\mathbb{E}[\hat{\boldsymbol{\beta}}_{\text{ridge}}] = (\mathbf{I}_p + \lambda (\mathbf{X}^T\mathbf{X})^{-1})^{-1} \boldsymbol{\beta} \neq \boldsymbol{\beta}$$
    $\Rightarrow$ Ridge estimates are \textbf{biased}.
    \item If predictors are standardized and uncorrelated:
    $$\hat{\beta}_{j,\text{ridge}} = \frac{1}{1 + \lambda} \hat{\beta}_{j,\text{OLS}}$$
    \item Variance of Ridge estimates is \textbf{smaller} than that of OLS, especially when data has multicollinearity.
\end{itemize}

\subsection{The Tuning Parameter $\lambda$}
\begin{itemize}
    \item $\lambda = 0$: Ridge reduces to OLS.
    \item $\lambda \to \infty$: All coefficients shrink to zero.
    \item $\lambda$ is chosen using \textbf{cross-validation} to minimize prediction error.
\end{itemize}

\subsection{Effective Degrees of Freedom}
In Ridge regression, the degrees of freedom is no longer simply $p$:
$$\text{df}(\lambda) = \text{tr}\left[(\mathbf{X}^T\mathbf{X} + \lambda\mathbf{I})^{-1}\mathbf{X}^T\mathbf{X}\right] = \sum_{j=1}^{p} \frac{d_j^2}{d_j^2 + \lambda}$$
where $d_j$ are the singular values of $\mathbf{X}$. As $\lambda \to \infty$, $\text{df}(\lambda) \to 0$.

\subsection{Why Ridge Keeps All Variables}
\begin{itemize}
    \item The $L_2$ penalty discourages large coefficients by shrinking them towards zero, but \textbf{never exactly to zero}.
    \item \textbf{Geometric intuition:} The constraint region for Ridge is a \textbf{ball} (circle/sphere). RSS contour ellipses typically touch the ball boundary at non-zero values.
    \item Hence, all predictors are \textbf{retained} in the model with reduced magnitudes.
\end{itemize}

\newpage
\section{LASSO Regression (L1 Regularization)}

\subsection{LASSO Formulation}
LASSO (Least Absolute Shrinkage and Selection Operator) minimizes:
$$\hat{\boldsymbol{\beta}}_{\text{lasso}} = \arg\min_{\boldsymbol{\beta}} \left\{ \sum_{i=1}^{n} \left( y_i - \beta_0 - \sum_{j=1}^{p} \beta_j x_{ij} \right)^2 + \lambda \sum_{j=1}^{p} |\beta_j| \right\}$$

\noindent\textbf{Equivalent constrained form:}
$$\min_{\boldsymbol{\beta}} \sum_{i=1}^{n} \left( y_i - \beta_0 - \sum_{j=1}^{p} \beta_j x_{ij} \right)^2 \quad \text{subject to} \quad \sum_{j=1}^{p} |\beta_j| \le t$$

\noindent\textbf{Matrix notation:}
$$\hat{\boldsymbol{\beta}}_{\text{lasso}} = \arg\min_{\boldsymbol{\beta}} \left\{ \frac{1}{2} \|\mathbf{y} - \mathbf{X}\boldsymbol{\beta}\|^2 + \lambda \|\boldsymbol{\beta}\|_1 \right\}$$

\subsection{Properties of LASSO Estimates}
\begin{itemize}
    \item \textbf{No closed-form solution:} The $L_1$ penalty involves absolute values, which are non-differentiable at zero. Solutions are nonlinear in $y_i$ due to the non-smooth constraint.
    \item \textbf{Biased} (towards zero), but with \textbf{smaller variance} than OLS.
    \item Does \textbf{not perform as well} as Ridge when data have strong multicollinearity problems (may unstably select only one variable from a correlated group).
    \item \textbf{Greatest advantage --- Sparsity:} LASSO can set some coefficients \textbf{exactly equal to zero} when $\lambda$ is sufficiently large, performing \textbf{automatic variable selection}.
\end{itemize}

\subsection{Sparsity: The Key Advantage}
In a model $Y = \beta_0 + \beta_1 X_1 + \cdots + \beta_p X_p + \epsilon$, we may believe many of the $\beta_j$'s are actually zero. LASSO seeks a set of \textbf{sparse solutions}:
\begin{itemize}
    \item \textbf{Geometric intuition:} The constraint region for LASSO is a \textbf{diamond} (in 2D) or a cross-polytope. RSS contour ellipses are more likely to touch the diamond at a \textbf{corner} (axis), where one or more coefficients are exactly zero.
    \item This means LASSO performs \textbf{model selection} automatically.
\end{itemize}

\subsection{Choosing $\lambda$ via Cross-Validation}
\begin{enumerate}
    \item Try a range of $\lambda$ values.
    \item For each $\lambda$: use the training set to train the model, predict values on the test/validation set.
    \item Compute the Root Mean Square Error:
    $$\text{RMSE} = \sqrt{\frac{\sum_{\text{test}} (y_i - \hat{y}_i)^2}{n_{\text{test}}}}$$
    \item Choose the $\lambda$ that has the \textbf{smallest RMSE} (or use $k$-fold CV and average).
\end{enumerate}

\newpage
\section{Comparison: OLS vs. Ridge vs. LASSO}

\begin{table}[h]
\centering
\small
\begin{tabularx}{\textwidth}{>{\raggedright\arraybackslash}p{0.20\textwidth}
                            >{\raggedright\arraybackslash}p{0.25\textwidth}
                            >{\raggedright\arraybackslash}p{0.25\textwidth}
                            >{\raggedright\arraybackslash}X}
\toprule
\textbf{Property} & \textbf{OLS} & \textbf{Ridge (L2)} & \textbf{LASSO (L1)} \\
\midrule
Objective & $\min \|\mathbf{y} - \mathbf{X}\boldsymbol{\beta}\|^2$ & $+ \lambda \sum \beta_j^2$ & $+ \lambda \sum |\beta_j|$ \\
Constraint shape & None & Ball (circle) & Diamond \\
Bias & Unbiased & Biased & Biased \\
Variance & Highest & Reduced & Reduced \\
Variable selection & No & No (all kept) & Yes (sparsity) \\
Closed-form & Yes & Yes & No \\
Best when & $p \ll n$, low collinearity & All predictors relevant, multicollinearity & Sparse true model \\
\bottomrule
\end{tabularx}
\end{table}

\end{document}
